% ============================================================================
%  SECCIÓN 1.4: FORMULACIÓN DEL PROBLEMA
% ============================================================================

\section{Formulaci\'on del problema}

La forma m\'as habitual de plantear un problema de investigaci\'on es mediante
una interrogaci\'on, aunque tambi\'en puede expresarse como un enunciado que
describa exactamente cu\'al es el problema que se pretende resolver. Para su
formulaci\'on se han considerado la magnitud, la trascendencia y la viabilidad
del problema, as\'i como su car\'acter factible y contrastable
\citep{arias2012}.

\subsection{Enunciado del problema}

El problema identificado se puede enunciar de la siguiente manera:

\begin{quote}
\textit{Insuficiente acceso de la ciudadan\'ia a informaci\'on confiable,
actualizada y georreferenciada sobre los escenarios deportivos del pa\'is, lo
que limita la pr\'actica deportiva y el aprovechamiento de la
infraestructura deportiva nacional.}
\end{quote}

El problema es factible de resolver porque la informaci\'on necesaria existe y
se genera de forma continua; lo que falta es un medio digital, centralizado e
interactivo que la organice y la ponga a disposici\'on de los usuarios. Es
contrastable, pues los resultados de la soluci\'on podr\'an verificarse mediante
indicadores de uso, cobertura y satisfacci\'on de los usuarios.

\subsection{Pregunta de investigaci\'on}

A partir del an\'alisis anterior, el problema se formula mediante la siguiente
pregunta de investigaci\'on:

\begin{quote}
\textit{¿C\'omo podr\'ia un sistema web y una aplicaci\'on m\'ovil contribuir a la
visualizaci\'on interactiva de los escenarios deportivos a nivel nacional,
facilitando a la ciudadan\'ia el acceso a informaci\'on confiable y actualizada
sobre su ubicaci\'on, sus caracter\'isticas y su oferta deportiva?}
\end{quote}

La respuesta a esta pregunta orienta el planteamiento de los objetivos de la
investigaci\'on, que constituyen el punto de referencia del estudio y se
desarrollan en la siguiente secci\'on.
